\chapter{Fisher Geometry and Invariants}
\label{ch:infogeometry}

Which geometric statements about the Gaussian recognition family are properties
of its laws, and which are properties of a chart or of a frame? This chapter
gives two answers. The first concerns the Fisher metric of the family. Its mean
sector equals the precision \(\Lambda\) in the moment chart, and in the
expectation chart it equals \(\Lambda\) on the quotient by the second-moment
directions but not on the coordinate mean directions, where it exceeds
\(\Lambda\) by a term governed by the squared Mahalanobis norm of the mean
(\Cref{prop:ig-fisher-expectation-chart}). The second concerns the interaction
precision. The generalized spectrum of the pencil \((L,\Lambda)\) is invariant
under every block reframing whenever the pencil is regular
(\Cref{prop:grg-pencil-invariance}), whereas ordinary eigenvalues and single
subspace determinants are frame dependent as soon as the block reframing group
contains the positive scalars (\Cref{prop:ig-frame-dependent-spectra-determinants}).
The pencil results type the singular boundary of the interaction family in
\Cref{ch:gaussian-renormalization}, and the invariance of agent-aligned law
functionals is what makes the determinant gap of
\Cref{thm:restrict-determinant-gap} a quantity about the model. The Fisher
metric is the tensor adopted in \Cref{ch:relational-inference}; its Hessian and
dual-inverse forms on the exponential charts are
\eqref{eq:ig-fisher-hessian} and \eqref{eq:ig-dual-inverse}, and the four charts
and the Frobenius pairing of symmetric-matrix coordinates are fixed in
\Cref{sec:ig-expfam}.

The chapter works with the Gaussian recognition family on the assembled latent
space \(\R^n\), \(n=NK\), in information parameters
\((\mathsf h,\Lambda)\) with \(\Lambda\succ0\), moments
\(\mu=\Lambda^{-1}\mathsf h\) and \(\Sigma=\Lambda^{-1}\), and second moment
\(M_2=\Sigma+\mu\mu^\top\). It does not assert that a rescaled coarse-graining
step contracts an information-geometric quantity along a history on which the
precision changes, that the dissipation identity of the fixed-covariance
relaxation of \Cref{sec:cg-gaussian-fixed-relaxation} is derived from the
fixed-precision pushforward of \Cref{prop:ig-galerkin-pushforward} or
conversely, or that an update other than that fixed-covariance mean flow is a
natural-gradient flow. \status{NOT-CLAIMED} The monotonicity question beyond a
fixed precision is \Cref{open:grg-galerkin-monotonicity}.

\section{The chart correction}
\label{sec:ig-chart}

It is often said that the mean sector of the Gaussian Fisher metric is the
precision. The statement is true in the moment chart \((\mu,\Sigma)\). In the
expectation chart \((\mu,M_2)\) it is true for the metric induced on the
quotient by the second-moment directions and false for the restriction to the
coordinate mean directions. The two propositions below compute both readings.

\propositionheading{Fisher metric in the moment chart}{prop:ig-fisher-moment-chart}
In the moment chart \((\mu,\Sigma)\), with \(\Lambda=\Sigma^{-1}\), the score of
the Gaussian family at \(y\) in the direction of a tangent vector
\((\dot\mu,\dot\Sigma)\), with \(\dot\Sigma\) symmetric, is
\begin{equation}
\partial_{(\dot\mu,\dot\Sigma)}\log q(y)
=\dot\mu^\top\Lambda(y-\mu)
+\tfrac12(y-\mu)^\top\Lambda\dot\Sigma\Lambda(y-\mu)
-\tfrac12\Tr(\Lambda\dot\Sigma),
\label{eq:ig-score}
\end{equation}
and the Fisher metric is block diagonal,
\begin{equation}
g^F_{(\mu,\Sigma)}\bigl((\dot\mu_1,\dot\Sigma_1),(\dot\mu_2,\dot\Sigma_2)\bigr)
=\dot\mu_1^\top\Lambda\dot\mu_2
+\tfrac12\Tr\left(\Lambda\dot\Sigma_1\Lambda\dot\Sigma_2\right).
\label{eq:ig-moment-metric}
\end{equation}
Every affine map \((\mu,\Sigma)\mapsto(Q\mu+\mathrm v,Q\Sigma Q^\top)\) with
\(Q\in\GL(n)\) and \(\mathrm v\in\R^n\) is an isometry of \eqref{eq:ig-moment-metric}.
\status{ESTABLISHED}

\noindent\emph{Proof.}
Write \(e=y-\mu\), so that
\(\log q(y)=-\tfrac12e^\top\Lambda e+\tfrac12\log\det\Lambda-\tfrac n2\log2\pi\).
A displacement \((\dot\mu,\dot\Sigma)\) moves \(e\) by \(-\dot\mu\) and
\(\Lambda\) by \(-\Lambda\dot\Sigma\Lambda\), and differentiating gives
\eqref{eq:ig-score}. Under \(e\sim\mathcal N(0,\Sigma)\) the linear term has
second moment \(\dot\mu^\top\Lambda\Sigma\Lambda\dot\mu=\dot\mu^\top\Lambda\dot\mu\),
the cross moment of the linear and quadratic terms is a third central moment of a
centered Gaussian and vanishes, and the Gaussian variance of a quadratic form
gives the quadratic term the variance
\(\tfrac14\operatorname{Var}(e^\top\Lambda\dot\Sigma\Lambda e)
=\tfrac12\Tr(\Lambda\dot\Sigma\Lambda\Sigma\Lambda\dot\Sigma\Lambda\Sigma)
=\tfrac12\Tr(\Lambda\dot\Sigma\Lambda\dot\Sigma)\). Polarization gives
\eqref{eq:ig-moment-metric}. Under the affine map the tangent vector becomes
\((Q\dot\mu,Q\dot\Sigma Q^\top)\) and the precision becomes
\(Q^{-\top}\Lambda Q^{-1}\), and both terms of \eqref{eq:ig-moment-metric} are
unchanged. \hfill$\square$

\medskip
This is the Fisher metric of the multivariate normal model whose Riemannian
geometry is developed by \citet{Skovgaard1984}; the form for a general
parameterization and the affine isometry are Eqs.~(7) and~(8) of
\citet{PineleStrapassonCosta2020}.

\propositionheading{Fisher metric in the expectation chart}{prop:ig-fisher-expectation-chart}
In the expectation chart \((\mu,M_2)\) the Fisher metric has the following
properties.
(i) Its mean block, cross block, and second block are
\begin{equation}
g_{[\mu\mu]}=\left(1+\mu^\top\Lambda\mu\right)\Lambda+\Lambda\mu\mu^\top\Lambda,
\label{eq:ig-meanblock}
\end{equation}
\begin{equation}
g_{[\mu M_2]}\bigl(\dot\mu,\dot M_2\bigr)=-\dot\mu^\top\Lambda\dot M_2\Lambda\mu,
\label{eq:ig-crossblock}
\end{equation}
and \(g_{[M_2M_2]}(\dot M_2,\dot M_2')=\tfrac12\Tr(\Lambda\dot M_2\Lambda\dot M_2')\),
so the metric is block diagonal in this chart exactly when \(\mu=0\).
(ii)\label{cor:ig-expectation-mean-quotient} The metric induced on the quotient
of the tangent space by the second-moment directions, equivalently the
restriction of \(g^F\) to the Fisher-orthogonal complement of those directions
read in the coordinate \(\mu\), is the Schur complement
\begin{equation}
g_{\mu}^{\mathrm{quot}}
=g_{[\mu\mu]}-g_{[\mu M_2]}g_{[M_2M_2]}^{-1}g_{[M_2\mu]}
=\Lambda.
\label{eq:ig-expectation-mean-quotient}
\end{equation}
(iii)\label{cor:ig-mean-block-discrepancy} With \(\hat\mu=\Lambda^{1/2}\mu\),
\begin{equation}
g_{[\mu\mu]}=\Lambda^{1/2}\left[\left(1+\lVert\hat\mu\rVert^2\right)I_n+\hat\mu\hat\mu^\top\right]\Lambda^{1/2},
\label{eq:ig-meanblock-factored}
\end{equation}
so the eigenvalues of \(\Lambda^{-1}g_{[\mu\mu]}\) are \(1+\mu^\top\Lambda\mu\)
with multiplicity \(n-1\) and \(1+2\mu^\top\Lambda\mu\) with multiplicity one;
hence \(g_{[\mu\mu]}\succeq\Lambda\), with equality if and only if \(\mu=0\).
\status{ESTABLISHED}

\noindent\emph{Proof.}
(i) The chart change is \(\Sigma=M_2-\mu\mu^\top\), so a displacement
\((\dot\mu,\dot M_2)\) induces
\begin{equation}
\delta\mu=\dot\mu,
\qquad
\delta\Sigma=\dot M_2-\left(\dot\mu\mu^\top+\mu\dot\mu^\top\right),
\label{eq:ig-chartjacobian}
\end{equation}
and the metric in this chart is the pullback of \eqref{eq:ig-moment-metric}
along \eqref{eq:ig-chartjacobian}. For the mean block put \(\dot M_2=0\); cycling
each of the four traces gives
\begin{equation}
\Tr\left(\Lambda\delta\Sigma\Lambda\delta\Sigma\right)
=2\left(\dot\mu^\top\Lambda\mu\right)^2+2\left(\mu^\top\Lambda\mu\right)\left(\dot\mu^\top\Lambda\dot\mu\right),
\label{eq:ig-crossterm}
\end{equation}
and adding the mean term of \eqref{eq:ig-moment-metric} to one half of
\eqref{eq:ig-crossterm} gives
\begin{equation}
\dot\mu^\top\Lambda\dot\mu+\left(\mu^\top\Lambda\mu\right)\left(\dot\mu^\top\Lambda\dot\mu\right)+\left(\dot\mu^\top\Lambda\mu\right)^2
=\dot\mu^\top\left[\left(1+\mu^\top\Lambda\mu\right)\Lambda+\Lambda\mu\mu^\top\Lambda\right]\dot\mu,
\label{eq:ig-meanquad}
\end{equation}
which is \eqref{eq:ig-meanblock}. Pairing \((\dot\mu,0)\) with \((0,\dot M_2)\), only
the covariance parts meet, and they give
\(\tfrac12\Tr\bigl(\Lambda[-(\dot\mu\mu^\top+\mu\dot\mu^\top)]\Lambda\dot M_2\bigr)
=-\dot\mu^\top\Lambda\dot M_2\Lambda\mu\). For \(\dot\mu=0\) nothing changes. The
cross block vanishes for all arguments exactly when \(\Lambda^2\mu=0\), that is when
\(\mu=0\). (ii) The metric is positive definite on this minimal family, so the
second block and its Schur complement are invertible, and the block-inverse
identity gives \((g^{\mathrm{quot}}_\mu)^{-1}=(g_{\boldsymbol\mu}^{-1})_{[\mathsf h\mathsf h]}\),
where \(g_{\boldsymbol\mu}\) is the metric in the expectation chart. By
\eqref{eq:ig-dual-inverse}, \(g_{\boldsymbol\mu}^{-1}\) is the metric in the natural chart, whose
block in the coordinate \(\mathsf h\) paired with \(\mu\) is
\(\operatorname{Cov}_q(Y)=\Sigma\) by \eqref{eq:ig-fisher-hessian}; inverting gives
\eqref{eq:ig-expectation-mean-quotient}. The Schur complement is the metric on the
orthogonal quotient. (iii) Equation \eqref{eq:ig-meanblock-factored} is
\eqref{eq:ig-meanblock} rewritten, and \(\Lambda^{-1}g_{[\mu\mu]}\) is similar to the
bracket, whose eigenvalues are \(1+2\lVert\hat\mu\rVert^2\) along \(\hat\mu\) and
\(1+\lVert\hat\mu\rVert^2\) on its orthogonal complement. Finally
\(g_{[\mu\mu]}-\Lambda=(\mu^\top\Lambda\mu)\Lambda+\Lambda\mu\mu^\top\Lambda\) is a sum
of two positive semidefinite matrices, which vanishes only if
\(\mu^\top\Lambda\mu=0\), that is \(\mu=0\). \hfill$\square$

\medskip
The two restriction mean blocks, \(\Lambda\) in \eqref{eq:ig-moment-metric} and
\eqref{eq:ig-meanblock}, differ because the coordinate vector
\(\partial/\partial\mu\) at fixed \(\Sigma\) and the one at fixed \(M_2\) are
different tangent vectors at the same law: the first moves the mean and changes
\(M_2\), the second moves the mean and compensates in \(\Sigma\). Their
discrepancy is governed by \(\mu^\top\Lambda\mu\) and is unbounded, while the
quotient metric equals \(\Lambda\) at every point. A coordinate identity that
mixes the mean and covariance sectors of a Gaussian family must therefore name
its chart and say whether it restricts or quotients.

\section{The mean sector and the interaction Laplacian}

The recognition family is a location family in its mean, so the mean parameter
takes values in the latent space \(\R^n\) itself, and the mean block of
\eqref{eq:ig-moment-metric} has the same components as the sample-space quadratic
form \(y\mapsto\tfrac12y^\top\Lambda y\). The coincidence is one of components in
matched charts and not of objects: one side is a metric on a manifold of laws,
the other a bilinear form on the latent space, and for a family that is not a
location family the two do not even have the same dimension. The next proposition
records the scoped relation to the Laplacian part \(L\) of the interaction family
of \Cref{def:gauss-interaction-family}.

\propositionheading{Edge relative entropy at fixed covariance}{prop:ig-fixed-covariance-kl}
Let the edge \((i,j)\) carry a positive-definite weight
\(W_{ij}\in\Sym_{++}^K\). For endpoint laws with means \(\mu_i,\mu_j\in\R^K\) and the
common covariance \(W_{ij}^{-1}\),
\begin{equation}
\KL\bigl(\mathcal N(\mu_i,W_{ij}^{-1})\Vert\mathcal N(\mu_j,W_{ij}^{-1})\bigr)
=\tfrac12(\mu_i-\mu_j)^\top W_{ij}(\mu_i-\mu_j)
=\tfrac12g^F_{(\mu,W_{ij}^{-1})}\bigl((\mu_i-\mu_j,0),(\mu_i-\mu_j,0)\bigr).
\label{eq:ig-fixed-covariance-kl-fisher}
\end{equation}
If every nonzero edge weight is positive definite, the sum of these relative
entropies over the edges equals \(\tfrac12\mu^\top L\mu\), the Laplacian part of the
energy \eqref{eq:gauss-interaction-energy} evaluated at the stacked mean \(\mu\).
\status{ESTABLISHED}

\noindent\emph{Proof.}
The relative entropy of two Gaussians with a common covariance \(\Sigma_0\) is
\(\tfrac12(\mu_i-\mu_j)^\top\Sigma_0^{-1}(\mu_i-\mu_j)\), and the mean block of
\eqref{eq:ig-moment-metric} at covariance \(W_{ij}^{-1}\) is \(W_{ij}\). The sum over
edges is the difference form
\(\sum_{i<j}(\mu_i-\mu_j)^\top W_{ij}(\mu_i-\mu_j)=\mu^\top L\mu\) of the Laplacian
part. \hfill$\square$

\medskip
The assembled \(NK\times NK\) matrix \(L\) enters only through this sum over
typed edge quadratics, and a singular weight has no fixed-covariance reading along
its kernel. For a regular non-Gaussian family the analogous quadratic is only the
second-order term of the relative entropy at a frozen base law
(\Cref{prop:pb-kl-divergence-jets}), so outside these Gaussian or frozen-local
hypotheses a sample-space Laplacian is neither a recognition Fisher tensor nor a
natural-gradient operator.

\section{Pencils and gauge invariants}
\label{sec:ig-gauge}
\label{sec:grg-spectra}

A change of local trivialization acts on every fiber and changes the components
of every object while changing no law. On the assembled latent coordinate it acts
agent by agent.

\definitionheading{Block reframing and the interaction pencil}{def:ig-block-reframing}
Let \(\Gamma_b=\rho_b(G)\) be the represented frame group of
\Cref{hyp:gen-lg-frame-group}. A block reframing is the invertible new-to-old
coordinate map
\begin{equation}
\mathbf P=\operatorname{diag}(\mathbf P_1,\dots,\mathbf P_N),
\qquad
\mathbf P_i\in\Gamma_b,
\qquad
y_{\mathrm{old}}=Py_{\mathrm{new}}.
\label{eq:ig-reframing}
\end{equation}
Split the interaction precision once, in the chart in which
\Cref{def:gauss-interaction-family} is imposed, as \(\Lambda=L+A\), with
\(L\) its Laplacian part and \(A=\operatorname{diag}(A_1,\dots,A_N)\) its self
terms. A block reframing acts on these quadratic forms by congruence,
\begin{equation}
\Lambda\longmapsto \mathbf P^\top\Lambda \mathbf P,
\qquad
L\longmapsto \mathbf P^\top L\mathbf P,
\qquad
A\longmapsto \mathbf P^\top A\mathbf P.
\label{eq:ig-congruence}
\end{equation}
The interaction pencil is the pair \((L,\Lambda)\); its generalized eigenvalues
are the roots \(\mathsf d\) of \(\det(L-\mathsf d\Lambda)\), with eigenvectors
\begin{equation}
Lx=\mathsf d\Lambda x.
\label{eq:ig-pencil}
\end{equation}

Because the reframing is block diagonal, \(\mathbf P^\top A\mathbf P\) is again block diagonal,
while \(\mathbf P^\top L\mathbf P\) need not be a matrix-weighted Laplacian when neighboring blocks
are reframed differently, since the row sums that define that normal form are not
preserved. The invariants below compare forms and are unaffected by that change of
normal form. The law \(\mathcal N(\mu,\Sigma)\) in old coordinates is
\(\mathcal N(\mathbf P^{-1}\mu,\mathbf P^{-1}\Sigma \mathbf P^{-\top})\) in new coordinates, so by the
affine isometry of \Cref{prop:ig-fisher-moment-chart} with \(Q=\mathbf P^{-1}\) every block
reframing is an isometry of the Fisher metric.

\propositionheading{Regular-pencil invariance}{prop:grg-pencil-invariance}
\label{prop:ig-generalized-spectrum-invariance}
Let \(L\) and \(\Lambda\) be symmetric with \(\det(L-\mathsf d\Lambda)\) not the zero
polynomial, and let \(\mathbf P\) be invertible. Then
\begin{equation}
\det\left(\mathbf P^\top L\mathbf P-\mathsf d\mathbf P^\top\Lambda \mathbf P\right)
=\det(\mathbf P)^2\det(L-\mathsf d\Lambda),
\label{eq:grg-pencil-congruence}
\end{equation}
so \((L,\Lambda)\) and \((\mathbf P^\top L\mathbf P,\mathbf P^\top\Lambda \mathbf P)\) have the same generalized
eigenvalues with the same multiplicities, and \(x\) solves \eqref{eq:ig-pencil}
exactly when \(\mathbf P^{-1}x\) solves the reframed problem with the same \(\mathsf d\). If
\(\Lambda\) is invertible, there are exactly \(n=NK\) finite roots.
\status{ESTABLISHED}

\noindent\emph{Proof.}
Equation \eqref{eq:grg-pencil-congruence} is multiplicativity of the determinant,
and the nonzero factor \(\det(\mathbf P)^2\) changes neither the roots nor their
multiplicities. If \(Lx=\mathsf d\Lambda x\), then
\(\mathbf P^\top L\mathbf P(\mathbf P^{-1}x)=\mathbf P^\top Lx=\mathsf d\mathbf P^\top\Lambda \mathbf P(\mathbf P^{-1}x)\), and the argument
with \(\mathbf P^{-1}\) gives the converse. When \(\Lambda\) is invertible the coefficient of
\(\mathsf d^{n}\) in \(\det(L-\mathsf d\Lambda)\) is \((-1)^n\det\Lambda\neq0\).
\hfill$\square$

\propositionheading{Regularity criterion}{prop:grg-pencil-regularity}
For \(L,A\succeq0\) and \(\Lambda=L+A\), the pencil \((L,\Lambda)\) is regular if and
only if
\begin{equation}
\ker L\cap\ker A=\{0\}.
\label{eq:grg-regularity}
\end{equation}
In that case \(\Lambda\succ0\). \status{ESTABLISHED}

\noindent\emph{Proof.}
Positive semidefiniteness gives \(\ker(L+A)=\ker L\cap\ker A\). A nonzero vector in
the intersection annihilates \(L-\mathsf d\Lambda\) for every \(\mathsf d\), so the
determinant vanishes identically. Conversely, if the intersection is trivial then
\(L+\Lambda=2L+A\succ0\), so the determinant is nonzero at \(\mathsf d=-1\), and
\(\ker\Lambda=\{0\}\). \hfill$\square$

\medskip
Regularity is load bearing. Congruence covariance holds for every pencil, but for
an irregular one the determinant vanishes identically. At the boundary operator
\(A=0\) of the interaction family one has \(\Lambda=L\), and
\(\det(L-\mathsf dL)\equiv0\) whenever \(\ker L\neq\{0\}\); restricting to a
complement of \(\ker L\) leaves the constant spectrum \(\{1\}\), which carries no
information about the graph. A quotient, pin, or mass repair must be declared and
its pencil proved regular before the statements below apply to such a point.

\propositionheading{Localization and endpoints}{prop:ig-generalized-spectrum-localization}
Let \(L\succeq0\) and \(A\succeq0\) with \(\Lambda=L+A\succ0\). Every generalized
eigenvalue of \eqref{eq:ig-pencil} is real and lies in \([0,1]\); the eigenspace of
\(0\) is \(\ker L\), and the eigenspace of \(1\) is \(\ker A\). If \(L\) is the
Laplacian part of \Cref{def:gauss-interaction-family}, every configuration that is
constant across agents lies in \(\ker L\), so \(\dim\ker L\geq K\).
\status{ESTABLISHED}

\noindent\emph{Proof.}
Because \(\Lambda\succ0\), the pencil is equivalent to the symmetric matrix
\(\Lambda^{-1/2}L\Lambda^{-1/2}\), so its \(n\) generalized eigenvalues are real and
are the stationary values of the Rayleigh quotient
\(x^\top Lx/x^\top\Lambda x\). Since \(x^\top\Lambda x=x^\top Lx+x^\top Ax\) with both
terms nonnegative and the sum positive, the quotient lies in \([0,1]\).
\(Lx=0\) is \eqref{eq:ig-pencil} with \(\mathsf d=0\), and \(Lx=\Lambda x\) is \(Ax=0\).
For the Laplacian part, \(x^\top Lx=\sum_{i<j}(x_i-x_j)^\top W_{ij}(x_i-x_j)\)
vanishes on every \(x=\mathbf 1\otimes w\) with \(w\in\R^K\), and these configurations
form a \(K\)-dimensional space. \hfill$\square$

\medskip
The generalized eigenvalue has a frame-independent reading: \(\mathsf d_k\) is the
fraction of the precision along the \(k\)-th pencil direction that interaction
supplies rather than self terms. Consensus directions sit at \(\mathsf d=0\),
because interaction supplies no precision there, and a direction that no self term
constrains sits at \(\mathsf d=1\).

\propositionheading{Mass-pencil transfer}{prop:grg-mass-pencil-transfer}
Let \(L\succeq0\), \(\Lambda_{\mathrm{ref}}\succ0\), \(\mathsf a>0\), and
\(\Lambda=L+\mathsf a\Lambda_{\mathrm{ref}}\). The generalized eigenvalues
\(\mathsf d_{\mathrm{ref}}\) of \((L,\Lambda_{\mathrm{ref}})\) and \(\mathsf d\) of
\((L,\Lambda)\) correspond by
\begin{equation}
\mathsf d=\frac{\mathsf d_{\mathrm{ref}}}{\mathsf d_{\mathrm{ref}}+\mathsf a},
\qquad
\mathsf d_{\mathrm{ref}}=\frac{\mathsf a\mathsf d}{1-\mathsf d},
\label{eq:grg-mass-transfer}
\end{equation}
with the same eigenvectors and multiplicities, and the correspondence preserves
order. Consequently the counting functions
\(\mathsf n(t)=\#\{k:\mathsf d_k\leq t\}\) and
\(\mathsf n_{\mathrm{ref}}(s)=\#\{k:\mathsf d_{\mathrm{ref},k}\leq s\}\) satisfy
\begin{equation}
\mathsf n(t)=\mathsf n_{\mathrm{ref}}\left(\frac{\mathsf at}{1-t}\right),
\qquad 0\leq t<1,
\label{eq:grg-counting-transfer}
\end{equation}
and if \(\mathsf n_{\mathrm{ref}}(s)-\mathsf n_{\mathrm{ref}}(0)\sim\mathrm c_1s^\alpha\) as
\(s\downarrow0\), then \(\mathsf n(t)-\mathsf n(0)\sim\mathrm c_1\mathsf a^\alpha t^\alpha\) as
\(t\downarrow0\). \status{ESTABLISHED}

\noindent\emph{Proof.}
For \(\mathsf d\neq1\), the equation \(Lv=\mathsf d(L+\mathsf a\Lambda_{\mathrm{ref}})v\)
is equivalent to \(Lv=\bigl(\mathsf a\mathsf d/(1-\mathsf d)\bigr)\Lambda_{\mathrm{ref}}v\),
and \(\mathsf d=1\) would require \(\Lambda_{\mathrm{ref}}v=0\), which is impossible.
Since \(L\succeq0\), every \(\mathsf d_{\mathrm{ref}}\) is nonnegative, and the map
\(\mathsf d_{\mathrm{ref}}\mapsto\mathsf d_{\mathrm{ref}}/(\mathsf d_{\mathrm{ref}}+\mathsf a)\)
is increasing on \([0,\infty)\); this gives the order, the counting identity, and,
because \(\mathsf at/(1-t)=\mathsf at(1+O(t))\), the asymptotic consequence. Without
\(L\succeq0\) the order can fail: for \(L=\operatorname{diag}(-2,1,3)\),
\(\Lambda_{\mathrm{ref}}=I\), and \(\mathsf a=1\), the roots \(-2,1,3\) map to
\(2,\tfrac12,\tfrac34\). \hfill$\square$

\medskip
Mass removal and the small-eigenvalue limit do not commute,
\begin{equation}
\lim_{\mathsf a\downarrow0}\lim_{\mathsf d_{\mathrm{ref}}\downarrow0}
\frac{\mathsf d_{\mathrm{ref}}}{\mathsf d_{\mathrm{ref}}+\mathsf a}=0,
\qquad
\lim_{\mathsf d_{\mathrm{ref}}\downarrow0}\lim_{\mathsf a\downarrow0}
\frac{\mathsf d_{\mathrm{ref}}}{\mathsf d_{\mathrm{ref}}+\mathsf a}=1,
\label{eq:grg-noncommuting-mass}
\end{equation}
so a positive mass regularizes each finite problem but supplies no mass-independent
spectrum at the boundary operator.

\propositionheading{Frame dependence of absolute spectra and subspace determinants}{prop:ig-frame-dependent-spectra-determinants}
(i) If \(\Gamma_b\supseteq\R_{>0}I_K\), in particular under the specialization
\(\Gamma_b=\GL^+(K)\) of \Cref{hyp:gen-lg-frame-group}, then \(\mathbf P=\mathrm cI_n\) with \(\mathrm c>0\)
is a block reframing and \(\operatorname{spec}(\mathbf P^\top\Lambda \mathbf P)=\mathrm c^2\operatorname{spec}(\Lambda)\).
(ii) Let \(\mathcal U'\) be a subspace of the new coordinate space with orthonormal basis
\(B'\), let \(\mathcal U=\mathbf P\mathcal U'\) have orthonormal basis \(B\), and let \(\Lambda'=\mathbf P^\top\Lambda \mathbf P\).
Then
\begin{equation}
\log\det\left(B'^\top\Lambda'B'\right)
=\log\det\left(B^\top\Lambda B\right)+\log\det\left(B'^\top \mathbf P^\top \mathbf PB'\right),
\label{eq:ig-logdet-law}
\end{equation}
so a subspace log-determinant is frame independent if and only if the restriction
of \(\mathbf P\) to \(\mathcal U'\) has unit absolute determinant; every orthogonal \(\mathbf P\) satisfies
this condition.
(iii) For a partition of the agents into blocks \(I\), the determinant gap
\(\tfrac12\bigl[\sum_I\log\det\Lambda_{II}-\log\det\Lambda\bigr]\) of
\Cref{thm:restrict-determinant-gap} is invariant under every block reframing. A
partition of the coordinates that splits an agent does not have this property: under
the specialization \(\Gamma_b=\GL^+(2)\), for \(N=1\), \(K=2\),
\(\Lambda=\bigl(\begin{smallmatrix}2&1\\1&2\end{smallmatrix}\bigr)\), and the blocks
\(\{1\}\), \(\{2\}\), the reframing
\(\mathbf P=\bigl(\begin{smallmatrix}1&0\\1&1\end{smallmatrix}\bigr)\) changes the gap from
\(\tfrac12\log\tfrac43\) to \(\tfrac12\log4\). \status{ESTABLISHED}

\noindent\emph{Proof.}
(i) \(\mathbf P^\top\Lambda \mathbf P=\mathrm c^2\Lambda\). (ii) Since \(\mathbf P\mathcal U'=\mathcal U\), write \(\mathbf PB'=BS\)
with \(S\) invertible. Then \(B'^\top\Lambda'B'=S^\top(B^\top\Lambda B)S\) and
\(B'^\top \mathbf P^\top \mathbf PB'=S^\top S\), so
\(\log\det(B'^\top\Lambda'B')=\log\det(B^\top\Lambda B)+\log\det(S)^2\) and
\(\log\det(B'^\top \mathbf P^\top \mathbf PB')=\log\det(S)^2\), which gives the identity. (iii) For blocks
that are unions of agents, the \((I,I)\) block of \(\mathbf P^\top\Lambda \mathbf P\) is
\(\mathbf P_I^\top\Lambda_{II}\mathbf P_I\) with \(\mathbf P_I=\operatorname{diag}(\mathbf P_i:i\in I)\), so
\(\sum_I\log\det\Lambda_{II}\) shifts by \(\sum_I\log\det(\mathbf P_I)^2=\log\det(\mathbf P)^2\), which
is the shift of \(\log\det\Lambda\). In the witness,
\(\mathbf P^\top\Lambda \mathbf P=\bigl(\begin{smallmatrix}6&3\\3&2\end{smallmatrix}\bigr)\) and
\(\det \mathbf P=1\), and the two gaps are \(\tfrac12\log(2\cdot2/3)\) and
\(\tfrac12\log(6\cdot2/3)\). \hfill$\square$

\medskip
A reframing changes no law, so a quantity computed from components that is not
invariant under \eqref{eq:ig-congruence} for every \(\mathbf P\) with blocks in \(\Gamma_b\) is
a property of the chart. When \(\Gamma_b\) contains the positive scalars, part (i)
lets a threshold on absolute eigenvalues of \(\Lambda\) accept or reject any given
direction by the choice of \(\mathrm c\), so such a threshold states nothing about the model.
Two invariant families remain available: generalized spectra of pencils whose two
members transform by the same congruence (\Cref{prop:grg-pencil-invariance}), and law
functionals of agent-aligned blocks, such as determinant gaps, mutual informations,
canonical correlations, and relative entropies between laws, which are invariant
because a reframing changes no law and which part (iii) exhibits in coordinates for
the gap. A construction that uses one subspace determinant must declare its frame,
carry the Jacobian term of \eqref{eq:ig-logdet-law}, or combine determinants so that
the Jacobians cancel, as the gap does. When \(\Gamma_b\subseteq\operatorname{O}(K)\), as
for a compact structure group, every block reframing is orthogonal and ordinary
spectra are invariant (\Cref{sec:gen-lg-frame-group}).

\propositionheading{Residual fiber label}{prop:grg-residual-fiber-label}
Assume \(\Gamma_b\supseteq\GL^+(K)\), as under the specialization of
\Cref{hyp:gen-lg-frame-group}. For \(M\in\PSD^K\), the orbit of \(M\) under global
reframing \(M\mapsto\bar{\mathbf P}^\top M\bar{\mathbf P}\), \(\bar{\mathbf P}\in\Gamma_b\), is the set of positive
semidefinite matrices of rank \(\operatorname{rank}M\). Hence the matrix factor \(M\)
of the biadditive family \eqref{eq:grg-biadditive} has \(K+1\) congruence classes and
no continuous modulus. \status{ESTABLISHED}

\noindent\emph{Proof.}
Congruence by an invertible matrix preserves rank and positive semidefiniteness. For
the converse, let \(r=\operatorname{rank}M\) and write \(M=QDQ^\top\) with \(Q\)
orthogonal and \(D\) diagonal, with positive entries \(d_1,\dots,d_r\) followed by
zeros. Then
\[
\bar{\mathbf P}=Q\operatorname{diag}\left(d_1^{-1/2},\dots,d_r^{-1/2},1,\dots,1\right)
\quad\text{gives}\quad
\bar{\mathbf P}^\top M\bar{\mathbf P}=\operatorname{diag}(I_r,0),
\]
and if \(\det\bar{\mathbf P}<0\), changing the
sign of the first column of \(\bar{\mathbf P}\) keeps this normal form and makes the
determinant positive. Two positive semidefinite matrices of equal rank are therefore
congruent to the same normal form by elements of \(\GL^+(K)\subseteq\Gamma_b\). This is
the positive semidefinite case of Sylvester's law of inertia
\citep[\S4.5]{HornJohnson2013}. \hfill$\square$

\medskip
The hypothesis on \(\Gamma_b\) is load bearing: when \(\Gamma_b\) is
\(\operatorname{SO}(K)\) or \(\operatorname{O}(K)\), the orbit of \(M\succ0\) is its
orthogonal-similarity class, labeled by \(K\) continuous eigenvalues
(\Cref{sec:gen-lg-frame-group}). The proposition also does not extend unchanged to a
nonflat pair \((W,\Theta)\), whose residual action combines congruence on \(W\) with
conjugation on \(\Theta\).
